% Bagian: computability
% Bab: computability-theory
% Subbagian: complete-ce-sets

\documentclass[../../../include/open-logic-section]{subfiles}

\begin{document}

\olfileid[id]{cmp}{thy}{cce}
\olsection{Himpunan Terenumerasi Secara Komputabel Lengkap}

\begin{defn}
Suatu himpunan $A$ merupakan \emph{himpunan !!{computably enumerable}
lengkap} (terhadap redusibilitas banyak-ke-satu) jika
\begin{enumerate}
\item $A$ terenumerasi secara komputabel, dan
\item untuk setiap himpunan terenumerasi secara komputabel lainnya $B$,
  $B \leq_m A$.
\end{enumerate}
\end{defn}

Dengan kata lain, himpunan terenumerasi secara komputabel lengkap adalah
himpunan terenumerasi secara komputabel yang ``paling sulit.'' Himpunan itu
memungkinkan kita menjawab pertanyaan tentang \emph{setiap} himpunan
terenumerasi secara komputabel.

\begin{thm}
$K$, $K_0$, dan $K_1$ semuanya merupakan himpunan terenumerasi secara
komputabel lengkap.
\end{thm}

\begin{proof}
Untuk melihat bahwa $K_0$ lengkap, misalkan $B$ sembarang himpunan yang
terenumerasi secara komputabel. Maka untuk suatu indeks $e$,
\[
B = W_e = \Setabs{x}{\cfind{e}(x) \fdefined}.
\]
Misalkan $f$ fungsi $f(x) = \tuple{e, x}$. Maka untuk setiap bilangan asli
$x$, $x \in B$ jika dan hanya jika $f(x) \in K_0$. Dengan kata lain, $f$
mereduksi $B$ ke~$K_0$.

Untuk melihat bahwa $K_1$ lengkap, perhatikan bahwa dalam bukti
\olref[k1]{prop:k1} kita mereduksi $K_0$ kepadanya. Jadi, berdasarkan
\olref[ppr]{prop:trans-red}, setiap himpunan terenumerasi secara komputabel juga
dapat direduksi ke~$K_1$.

Kelengkapan $K$ mengikuti dengan cara yang sama dari reduksi $K_0$ ke~$K$.
\end{proof}

\begin{prob}
Berikan reduksi $K$ ke $K_0$.
\end{prob}

\begin{digress}
Ternyata semua contoh himpunan terenumerasi secara komputabel yang telah kita
tinjau sejauh ini dapat dihitung atau lengkap. Hal ini seharusnya terasa
aneh!{} Adakah contoh himpunan terenumerasi secara komputabel yang tidak dapat
dihitung sekaligus tidak lengkap? Jawabannya ya, tetapi fakta ini baru
ditetapkan secara independen oleh Friedberg dan Muchnik pada pertengahan
dekade 1950-an.
\end{digress}

\end{document}

